\section{Scope and Motivation}

This methodological manuscript asks whether latent regulatory mechanisms can be identified from noisy, discretely sampled action trajectories. Attachment-related regulation provides the motivating application, but the methodological claim is broader than attachment theory. P, PI, and PID are treated as deliberately simple benchmark generators rather than privileged psychological mechanisms.

System identification requires more than fitting a plausible trajectory. A model family can predict well while its parameters or mechanism remain practically non-identifiable, especially under limited data, measurement error, feedback, nonlinear saturation, and partial observation \parencite{ljung1999systemidentification,lam2022identifiability}. This distinction is critical when model labels are interpreted psychologically.

The paper therefore separates three questions: whether a model predicts held-out behavior; whether the correct generating family can be recovered; and whether parameters within a selected family are estimated with useful precision. Failure on the second or third question limits mechanistic interpretation even if the first question is satisfactory.
